% ==============================================================================
% sci_s2_setup_sensors.tex — SCI/SICE tutorial Session 2:
% environment setup and sensor data acquisition
% SCI/SICE チュートリアル セッション2: 開発環境のセットアップとセンサデータの取得
% ==============================================================================

% ------------------------------------------------------------------
\begin{frame}{地図}
  \centering
  \vspace{-1mm}
  \resizebox{0.62\textwidth}{!}{\scisessionmap{2}}
\end{frame}

% ------------------------------------------------------------------
\begin{frame}{要点3つ}
  \begin{enumerate}
    \item ビルド環境は ESP-IDF と \code{sf CLI} が担う。開発を伴わない参加者は Web Flasher / GUI インストーラという代替経路がある
    \item 学習者が書くコードは \api{setup()} と \api{loop\_400Hz(dt)} の2関数だけ。センサやHWの詳細は \api{ws::} 名前空間が隠す
    \item センサデータの取得経路は3通りある。目的に応じて \api{ws::print}、\code{sf telemetry}、\code{sf log wifi} を使い分ける
  \end{enumerate}
\end{frame}

% ------------------------------------------------------------------
\begin{frame}[fragile]{開発環境の導入 (1/2): PC 側}
  \begin{block}{方法①: GUI インストーラ「StampFly Setup」}
    \footnotesize ターミナル操作なしで導入できる。GitHub の Releases ページから OS 別の実行ファイルをダウンロードして起動し、ウィザードに従うだけで \code{sf} CLI 一式（ESP-IDF: Espressif の ESP32 向け開発基盤を含む）が入る\\
    \url{https://github.com/M5Fly-kanazawa/stampfly_ecosystem/releases/latest}
  \end{block}

  \vspace{0.15em}

  \begin{exampleblock}{方法②: CLI インストーラ}
    \begin{lstlisting}[style=sfbash, numbers=none]
git clone https://github.com/M5Fly-kanazawa/stampfly_ecosystem
cd stampfly_ecosystem
./install.sh          # Windows は install.bat
    \end{lstlisting}
  \end{exampleblock}

  \vspace{0.15em}

  {\scriptsize いずれの方法で導入しても、この後の \code{source setup\_env.sh} $\to$ \code{sf doctor}（このあとの「事前準備の確認」で）確認する}
\end{frame}

% ------------------------------------------------------------------
\begin{frame}{開発環境の導入 (2/2): 機体とペアリング}
  \begin{block}{開発環境なしでも書き込みだけはできる}
    \small ブラウザ（Chrome/Edge）から Web Flasher を開き、公開済みの完成版ファームウェアを機体・コントローラへ直接書き込める。インストール不要\\
    \url{https://m5fly-kanazawa.github.io/stampfly_ecosystem/flash/}
  \end{block}

  \vspace{0.5em}

  \begin{exampleblock}{コントローラのペアリング（初回のみ）}
    \footnotesize
    \begin{enumerate}\setlength{\itemsep}{2pt}
      \item コントローラの LCD パネルボタンを押しながら電源を入れる
      \item StampFly 本体のボタンを2秒長押しする
      \item 双方がビープしたらペアリング完了。以降は電源を入れるだけで自動再接続する
    \end{enumerate}
  \end{exampleblock}
\end{frame}

% ------------------------------------------------------------------
\begin{frame}[fragile]{事前準備の確認}
  \begin{block}{全員で \code{sf doctor} を実行する}
    \small ESP-IDF、Python、USBシリアルドライバ、sf CLI 自体のバージョンをまとめて診断する
  \end{block}

  \begin{lstlisting}[style=sfbash, numbers=none]
source setup_env.sh
sf doctor
  \end{lstlisting}

  \vspace{0.3em}

  \begin{alertblock}{ここで通らなくても心配は不要}
    \footnotesize すべて OK にならなかった場合は、このあとのデモは「見るだけ」で参加し、復習パス（本セッション末尾）に沿って会場外で環境を整えてほしい
  \end{alertblock}
\end{frame}

% ------------------------------------------------------------------
\begin{frame}{ツールチェーン}
  \begin{block}{ESP-IDF の上に sf CLI を重ねる}
    \small \code{idf.py} を直接叩く代わりに、ビルド・書き込み・診断・ログ取得を統一したコマンド体系で提供する
  \end{block}

  \vspace{0.5em}

  \begin{table}
    \centering\small
    \begin{tabular}{ll}
      \hline
      \textbf{コマンド} & \textbf{役割} \\
      \hline
      \code{sf build [target]} & ファームウェアビルド \\
      \code{sf flash [target] -m} & 書き込み（\code{-m} でモニタ付き） \\
      \code{sf monitor}          & シリアルモニタ \\
      \code{sf doctor}           & 環境診断 \\
      \hline
    \end{tabular}
  \end{table}

  \vspace{0.5em}

  {\small \code{target} には \code{vehicle} / \code{controller} / \code{workshop} を指定する}
\end{frame}

% ------------------------------------------------------------------
\begin{frame}[fragile]{学習者コードの型とワークフロー}
  \begin{columns}[T]
    \begin{column}{0.44\textwidth}
      \begin{block}{学習者が書く関数は2つだけ}
        \footnotesize
        \begin{itemize}
          \item \api{setup()} --- 起動時に1回
          \item \api{loop\_400Hz(dt)} --- 400\,Hz毎
        \end{itemize}
      \end{block}
      \vspace{0.3em}
      \begin{exampleblock}{切替・ビルド・書き込み}
        \footnotesize
        \code{sf lesson switch sci2026:N}\\
        \code{sf lesson build}\\
        \code{sf lesson flash}
      \end{exampleblock}
    \end{column}
    \begin{column}{0.54\textwidth}
      \centering
      \resizebox{\columnwidth}{!}{\input{../../_shared/tikz/build_flash_flow}}
    \end{column}
  \end{columns}
\end{frame}

% ------------------------------------------------------------------
\begin{frame}{IMU の軸定義}
  \vspace{-0.5em}
  \begin{center}
    \resizebox{0.6\textwidth}{!}{\documentclass[border=5pt]{standalone}
\usepackage{tikz}
\usetikzlibrary{arrows.meta, calc, decorations.markings}

\begin{document}

\definecolor{axisX}{RGB}{220, 50, 30}   % Red - X forward (North)
\definecolor{axisY}{RGB}{40, 160, 40}   % Green - Y right (East)
\definecolor{axisZ}{RGB}{0, 120, 200}   % Blue - Z down
\definecolor{bodycolor}{RGB}{80, 80, 80}

\begin{tikzpicture}[
    axis/.style={-{Stealth[length=5mm, width=3mm]}, line width=2.5pt},
    rotarrow/.style={
      -{Stealth[length=3.5mm, width=2.5mm]},
      line width=1.8pt,
    },
  ]

  % Isometric unit vectors
  \pgfmathsetmacro{\Xx}{-0.71}
  \pgfmathsetmacro{\Xy}{0.50}
  \pgfmathsetmacro{\Yx}{1.0}
  \pgfmathsetmacro{\Yy}{0.0}
  \pgfmathsetmacro{\Zx}{0.0}
  \pgfmathsetmacro{\Zy}{-1.0}

  % Body - isometric parallelogram
  \fill[bodycolor!10]
    ({1.1*\Xx+1.1*\Yx},{1.1*\Xy+1.1*\Yy}) --
    ({1.1*\Xx-1.1*\Yx},{1.1*\Xy-1.1*\Yy}) --
    ({-1.1*\Xx-1.1*\Yx},{-1.1*\Xy-1.1*\Yy}) --
    ({-1.1*\Xx+1.1*\Yx},{-1.1*\Xy+1.1*\Yy}) -- cycle;
  \draw[bodycolor, line width=1.5pt]
    ({1.1*\Xx+1.1*\Yx},{1.1*\Xy+1.1*\Yy}) --
    ({1.1*\Xx-1.1*\Yx},{1.1*\Xy-1.1*\Yy}) --
    ({-1.1*\Xx-1.1*\Yx},{-1.1*\Xy-1.1*\Yy}) --
    ({-1.1*\Xx+1.1*\Yx},{-1.1*\Xy+1.1*\Yy}) -- cycle;

  % Axis length
  \pgfmathsetmacro{\AL}{4.5}

  % X axis (Forward / North)
  \draw[axis, axisX] (0,0) -- ({\AL*\Xx},{\AL*\Xy});

  % Front marker triangle -- drawn AFTER the X axis line so its opaque fill
  % paints over the line where it would otherwise cross the marker,
  % keeping the marker readable as one solid shape (its base midpoint sits
  % exactly on the X-axis line, which used to bisect it when drawn first).
  % 前方マーカーの三角形。X軸の線より後に描く（三角形の塗りが不透明なため、
  % そのままだと線が三角形を貫通する箇所を上塗りし、マーカーが1つの
  % 塗りつぶし形状として読めるようにする）。三角形の底辺の中点はちょうど
  % X軸の直線上にあり、先に描くと三角形が線で分断されて見えていた。
  \fill[axisX!50]
    ({1.1*\Xx-0.35},{1.1*\Xy}) --
    ({1.1*\Xx+0.35},{1.1*\Xy}) --
    ({1.1*\Xx},{1.1*\Xy+0.3}) -- cycle;

  \node[font=\Large\bfseries, axisX, above left]
    at ({\AL*\Xx},{\AL*\Xy}) {X (Forward)};

  % Y axis (Right / East)
  \draw[axis, axisY] (0,0) -- ({\AL*\Yx},{\AL*\Yy});
  \node[font=\Large\bfseries, axisY, above right]
    at ({\AL*\Yx},{\AL*\Yy}) {Y (Right)};

  % Z axis (Down)
  \draw[axis, axisZ] (0,0) -- ({\AL*\Zx},{\AL*\Zy});
  \node[font=\Large\bfseries, axisZ, right]
    at ({\AL*\Zx},{\AL*\Zy}) {Z (Down)};

  % === Rotation arrows (elliptical arcs wrapping around each axis) ===

  % Roll - rotation around X axis
  % Arc lies in the Y-Z plane, centered on X axis
  % Projected semi-axes: Y direction (1,0) and Z direction (0,-1)
  \pgfmathsetmacro{\Rcx}{2.6*\Xx}  % center on X axis
  \pgfmathsetmacro{\Rcy}{2.6*\Xy}
  \pgfmathsetmacro{\Rr}{0.9}        % arc radius
  \draw[rotarrow, axisX!70!black]
    ({\Rcx+\Rr*\Yx*cos(200)+\Rr*\Zx*sin(200)},
     {\Rcy+\Rr*\Yy*cos(200)+\Rr*\Zy*sin(200)})
    \foreach \a in {200,210,...,490} {
      -- ({\Rcx+\Rr*\Yx*cos(\a)+\Rr*\Zx*sin(\a)},
          {\Rcy+\Rr*\Yy*cos(\a)+\Rr*\Zy*sin(\a)})
    };
  \node[font=\normalsize\bfseries, axisX!70!black, above]
    at ({\Rcx},{\Rcy+\Rr+0.15}) {Roll};

  % Pitch - rotation around Y axis
  % Arc lies in the X-Z plane, centered on Y axis
  % Projected semi-axes: X direction (-0.71,0.5) and Z direction (0,-1)
  %
  % Right-hand cyclic rule for a positive rotation about axis A: the arc
  % must be parametrised as (first-of-remaining-pair)*cos + (second)*sin
  % following the cyclic order X->Y->Z->X. Roll (about X) uses Y*cos+Z*sin
  % and yaw (about Z) uses X*cos+Y*sin, so pitch (about Y) must use
  % Z*cos+X*sin -- NOT X*cos+Z*sin, which would reverse the sense.
  % 右手系の巡回則: 軸Aまわりの正回転の円弧は、X->Y->Z->Xの巡回順で
  % 「（残り2軸のうち先の軸）*cos + （後の軸）*sin」としてパラメータ化する。
  % ロール（X軸まわり）はY*cos+Z*sin、ヨー（Z軸まわり）はX*cos+Y*sinを
  % 使うので、ピッチ（Y軸まわり）はZ*cos+X*sinでなければならない
  % （X*cos+Z*sinでは回転の向きが逆になる）。
  \pgfmathsetmacro{\Pcx}{2.6*\Yx}
  \pgfmathsetmacro{\Pcy}{2.6*\Yy}
  \pgfmathsetmacro{\Pr}{0.9}
  \draw[rotarrow, axisY!70!black]
    ({\Pcx+\Pr*\Zx*cos(200)+\Pr*\Xx*sin(200)},
     {\Pcy+\Pr*\Zy*cos(200)+\Pr*\Xy*sin(200)})
    \foreach \a in {200,210,...,490} {
      -- ({\Pcx+\Pr*\Zx*cos(\a)+\Pr*\Xx*sin(\a)},
          {\Pcy+\Pr*\Zy*cos(\a)+\Pr*\Xy*sin(\a)})
    };
  \node[font=\normalsize\bfseries, axisY!70!black, above right]
    at ({\Pcx+\Pr*0.3},{\Pcy+\Pr*0.7}) {Pitch};

  % Yaw - rotation around Z axis
  % Arc lies in the X-Y plane, centered on Z axis
  % Projected semi-axes: X direction (-0.71,0.5) and Y direction (1,0)
  \pgfmathsetmacro{\Wcx}{2.6*\Zx}
  \pgfmathsetmacro{\Wcy}{2.6*\Zy}
  \pgfmathsetmacro{\Wr}{0.9}
  \draw[rotarrow, axisZ!70!black]
    ({\Wcx+\Wr*\Xx*cos(160)+\Wr*\Yx*sin(160)},
     {\Wcy+\Wr*\Xy*cos(160)+\Wr*\Yy*sin(160)})
    \foreach \a in {160,170,...,450} {
      -- ({\Wcx+\Wr*\Xx*cos(\a)+\Wr*\Yx*sin(\a)},
          {\Wcy+\Wr*\Xy*cos(\a)+\Wr*\Yy*sin(\a)})
    };
  % Label pushed further outside the loop than the 0.8/0.3 multipliers
  % originally used (found by visual review to have its "w" overlapped
  % by the arc stroke).
  % ラベルを元の0.8/0.3倍率より外側へ押し出す（目視レビューで"w"の文字に
  % 弧の線が重なっていたのを修正）。
  \node[font=\normalsize\bfseries, axisZ!70!black, above left]
    at ({\Wcx-\Wr*1.4},{\Wcy+\Wr*0.55}) {Yaw};

\end{tikzpicture}
\end{document}
}
  \end{center}
  \vspace{-0.5em}
  \begin{block}{}
    \centering\small
    BMI270 --- 6軸（3軸ジャイロ + 3軸加速度）400\,Hz \quad|\quad \mbox{Body FRD（前方-右-下）座標系}
  \end{block}
\end{frame}

% ------------------------------------------------------------------
\begin{frame}{IMU API と単位}
  \begin{table}
    \centering\small
    \begin{tabular}{lll}
      \hline
      \textbf{関数} & \textbf{説明} & \textbf{単位} \\
      \hline
      \api{gyro\_x/y/z()}  & 角速度（Roll/Pitch/Yaw） & rad/s \\
      \api{accel\_x/y/z()} & 加速度（X/Y/Z） & m/s\textsuperscript{2} \\
      \hline
    \end{tabular}
  \end{table}

  \vspace{0.5em}

  \begin{exampleblock}{静止時の目安}
    \small ジャイロ $\approx 0$、加速度 Z $\approx -9.81$\,m/s\textsuperscript{2}（重力に対する反力）
  \end{exampleblock}
\end{frame}

% ------------------------------------------------------------------
\begin{frame}[fragile]{データ取得経路① Teleplot}
  \begin{block}{\texttt{ws::print} をそのままグラフにする}
    \small \code{>変数名:値} の形式で出力すると、VSCode拡張機能 Teleplot がリアルタイムにグラフ化する
  \end{block}

  \begin{sflisting}[basicstyle=\ttfamily\scriptsize]{user\_code.cpp}
static uint32_t tick = 0;
void loop_400Hz(float dt) {
    if (tick++ % 4 == 0) {   // 100Hz decimation
        ws::print(">gyro_x:%.3f", ws::gyro_x());
        ws::print(">gyro_y:%.3f", ws::gyro_y());
    }
}
  \end{sflisting}

  \vspace{-4.5mm}

  {\footnotesize VSCode 拡張 \code{alexnesnes.teleplot} を導入後、\code{sf monitor workshop} で接続（\code{sf lesson flash} 直後はそのまま開く）}
\end{frame}

% ------------------------------------------------------------------
\begin{frame}[fragile]{データ取得経路② sf telemetry}
  \begin{block}{コード変更なしでライブ表示}
    \small 機体は状態推定値を常時 50\,Hz で送信している。受信して表示するだけなら学習者コードの変更は不要
  \end{block}

  \begin{lstlisting}[style=sfbash, numbers=none]
sf telemetry          # ターミナル表示（既定）
sf telemetry --web    # ブラウザ表示（UDP:5005 -> SSE）
  \end{lstlisting}

  \vspace{0.3em}

  {\small 講義や複数人での画面共有には \code{--web} が見やすい。UDP:5005 を SSE（サーバーからブラウザへの一方向配信）でブラウザへ届けている}
\end{frame}

% ------------------------------------------------------------------
\begin{frame}[fragile]{データ取得経路③ ログ取得と可視化}
  \begin{block}{WiFi 経由で 400\,Hz Data Stream を記録する}
    \small \code{sf log wifi} は機体と同じ WiFi（SoftAP モードなら機体が出す AP）で 400\,Hz の Data Stream を UDP 受信し、\code{-o *.csv} で 1 周期 1 行の CSV に保存する
  \end{block}

  \begin{lstlisting}[style=sfbash, numbers=none]
sf log wifi -d 30 -o flight.csv   # 30秒キャプチャ -> Data Stream CSV
sf log viz flight.csv
  \end{lstlisting}

  \vspace{0.3em}

  \begin{exampleblock}{使い分けの目安}
    \footnotesize Teleplot = その場で波形を見る／\code{sf telemetry} = コード変更なしで確認／\code{sf log wifi} = 後で解析するためのデータ保存
  \end{exampleblock}
\end{frame}

% ------------------------------------------------------------------
\begin{frame}[fragile]{デモ: 傾けて確認}
  \vspace{0.3em}

  {\small 見るだけ ｜ \textcolor{sfblue}{\textbf{一緒に打つ}} ｜ 帰宅後に再現}

  \vspace{0.5em}

  \begin{block}{何を見るか}
    \begin{itemize}\small
      \item StampFly を手に持ち、前後左右に傾ける
      \item \api{gyro\_x}/\api{gyro\_y} が傾ける速さに応じて振れる
      \item 静止させると \api{accel\_z} が $-9.81$ 付近に戻る
    \end{itemize}
  \end{block}

  \begin{alertblock}{手順}
    \footnotesize \code{sf lesson switch sci2026:2} $\to$ データ取得経路①（Teleplot）のページのコードを書く $\to$ \code{sf lesson build \&\& sf lesson flash} $\to$ Teleplot で確認
  \end{alertblock}
\end{frame}

% ------------------------------------------------------------------
\begin{frame}{他のセンサ}
  \small IMU 以外のセンサも同じ考え方で読める。時間の都合で本セッションでは扱わないが、単独ビルド可能な例題として用意してある

  \vspace{0.5em}

  \begin{table}
    \centering\small
    \begin{tabular}{ll}
      \hline
      \textbf{例題} & \textbf{センサ} \\
      \hline
      \code{examples/05\_read\_tof} & VL53L3CX（ToF距離センサ） \\
      \code{examples/06\_read\_baro} & BMP280（気圧センサ） \\
      \hline
    \end{tabular}
  \end{table}

  \vspace{0.5em}

  {\footnotesize \code{firmware/vehicle/examples/} 配下。各ディレクトリで \code{idf.py set-target esp32s3 \&\& idf.py build} すれば単体で動く}
\end{frame}

% ------------------------------------------------------------------
\begin{frame}{理論とコードの対応表}
  \begin{table}
    \centering\small
    \begin{tabular}{lccc}
      \hline
      \textbf{理論} & \textbf{コード} & \textbf{Data Stream 列} & \textbf{トピック} \\
      \hline
      角速度 $\omega$ [rad/s]、Body FRD & \api{ws::gyro\_x()} & \code{gyro\_x} & \code{sensor\_imu} \\
      並進加速度 $a$ [m/s\textsuperscript{2}] & \api{ws::accel\_x()} & \code{accel\_x} & \code{sensor\_imu} \\
      \hline
    \end{tabular}
  \end{table}

  \vspace{0.5em}

  {\small 学習者が呼ぶ \api{ws::} 関数は、内部で \code{sensor\_imu} トピックを購読しているだけである
  （アクセス階層 L0（Workshop API）$\to$ L1（Topic API）の対応）}
\end{frame}

% ------------------------------------------------------------------
\begin{frame}{復習パス}
  \resizebox{\textwidth}{!}{%
  \begin{tabular}{lll}
    \hline
    \textbf{コマンド} & \textbf{実習} & \textbf{文書} \\
    \hline
    \code{sf lesson switch sci2026:1} & 実習1（環境セットアップ） & \code{docs/guides/troubleshooting.md} \\
    \code{sf lesson switch sci2026:2} & 実習2（IMU） & \code{firmware/vehicle/docs/architecture.md \S2} \\
    \code{sf log wifi} / \code{sf log viz} & --- & \code{docs/guides/flight-log-viz.md} \\
    \hline
  \end{tabular}}
\end{frame}

% ------------------------------------------------------------------
\begin{frame}{チェックポイント}
  \begin{alertblock}{確認事項}
    \begin{itemize}
      \item[$\square$] \code{sf doctor} が通る、または代替経路（Web Flasher / GUI インストーラ）を把握した
      \item[$\square$] \api{setup()}/\api{loop\_400Hz(dt)} の2関数構造を説明できる
      \item[$\square$] ジャイロ・加速度をどれか1つの経路で確認できた
    \end{itemize}
  \end{alertblock}

  \vspace{1em}

  \begin{block}{次のセッション}
    S3: モータ制御とコントローラ入力の実装 --- センサの次は、機体を動かす側を作る
  \end{block}
\end{frame}
